\documentclass{standalone}
\usepackage{standalone}

\begin{document}
\subsection{Contour Analysis}
Canny edge detection works perfectly for most cases and provides an easy to detect contours around plate borders. After passing the canny image to the OpenCV's $findContours$ function, we get a set of contours. For each of these contours several parameters are checked before selecting a contours as a possible plate.

In this steps, some of the estimated candidate from previous stages fall out. In Figure \ref{fig:ContourImage}, all detected contours are highlighted.
\begin{figure}
    \centering
    \includegraphics[width=.7\linewidth]{./img/sample/stage9.jpg}
    \caption{Only remaining candidate after contour analysis}
    \label{fig:ContourImage}
\end{figure}



\subsubsection{Width and Height}
The width and height of the detected contour should be above a minimum margin to be recognizable by the OCR. We set the minimum width = $30$ pixels and minimum height = $75$ pixels for estimated plate.

\subsubsection{Area Ratio}
Next we check if the area of the contour is greater than the $10\%$ of the entire image. Remember, we got this image by localizing the plate like regions in pre-processing stage. A possible plate should not have the area below $10\%$ of the entire image.
\begin{equation}
area\_ratio = \dfrac{contour\_area}{image\_area }
\end{equation}

\subsubsection{Aspect Ratio}
Equation \ref{eq:AspectRatio} shows the calculation of aspect ratio of the contour image. From review and observations, we set the minimum value of aspect ratio to $0.3$ and maximum to $0.6$.
\begin{equation} \label{eq:AspectRatio}
aspect\_ratio = \dfrac{contour\_height}{contour\_width}
\end{equation}

\subsubsection{Rotation}
The rotation can be found by the $minAreaRect$ function from OpenCV. 
\begin{lstlisting}[language=Python]
angle = cv2.minAreaRect(cnt)[0]
\end{lstlisting}
If $10 \leq |angle| \leq 80$ we skipped the contour. Because otherwise the contour will be too much skewed to be recognized by any OCR.

\end{document}